\section{Model Library}
\label{sec:model_library}

The framework currently provides implementations of published PFDHA models for the four model categories defined in Section~\ref{sec:pfdha_framework}. Table~\ref{tab:model_library} summarises the inventory by PFDHA component and faulting style. The principal components are listed separately as principal surface-rupture probability and principal fault displacement. Distributed-faulting models are separated where the published formulation provides distinct rupture-probability and displacement-amplitude terms; coupled SR+FD pathways are grouped together when the published model treats them as a paired implementation. The reference set draws from the FDHI Project \citep{Bozorgnia2021, Sarmiento2025}, the SURE~2.0 database \citep{Nurminen2022, Visini2025}, IAEA benchmark exercises \citep{valentini_probabilistic_2021, Valentini2025}, and earlier foundational studies \citep{youngs_methodology_2003, Petersen2011, MossRoss2011, Takao2013}. The framework currently includes implementations of over 20 published models across the four categories.


\subsection{Organisation and Displacement-Metric Heterogeneity}
\label{sec:model_library:organisation}

Each implemented model is a self-contained subclass of one of the four model-component base classes introduced in Section~\ref{sec:architecture:adapter} and is made available by name so that it can be selected from the model-logic-tree NRML file. Where a published study provides paired models---for example, a principal surface-rupture probability model and a companion principal displacement model that share calibration data---each component is implemented as a separate class so it can be selected independently in the logic tree, provided that the resulting branch remains consistent with the model's applicable faulting style, displacement metric, and validity range.

Across the principal displacement category, models predict different displacement metrics: vertical, horizontal (lateral), or net displacement components, and either single-principal, sum-of-principal, or aggregate displacement definitions; \citet[Table~1]{Sarmiento2025} provides the formal taxonomy. Mixing these metrics within a logic tree is not automatically meaningful. For example, comparing a vertical single-principal estimate from \citet{youngs_methodology_2003} with a net aggregate estimate from \citet{Kuehn2024} requires explicit assumptions about the slip vector, rupture geometry, and participation of distributed deformation. The framework therefore exposes the metric directly: each entry in Table~\ref{tab:model_library} is annotated with its component (v, h, or n) and its single-principal (S) or sum-of-principal/aggregate (A) classification. It is the analyst's responsibility to assemble metric-consistent logic-tree branches.

The model inventory also defines the available branch choices for model-logic-tree calculations. Each branch selects one model for each active PFDHA component. Branch combinations must remain consistent with the applicable faulting style, displacement metric, model validity range, and database assumptions of the selected published models, and---when a source-model logic tree is also in use---must remain physically consistent with the source-model branch under which they are evaluated (for example, a reverse-faulting FDHA-model branch should not be paired with a SMLT branch in which a \texttt{simpleFaultDipAbsolute} or \texttt{setMSRAbsolute} modification has changed the source kinematics or scaling regime). The logic-tree workflow therefore provides a reproducible mechanism for exploring epistemic uncertainty in both the source characterization and PFDHA model selection, but it does not make all possible cross-category, cross-tree combinations scientifically interchangeable.


\subsection{Principal Surface-Rupture Probability}
\label{sec:model_library:principal_sr}

Principal surface-rupture probability models predict $P[\,sr \neq 0 \mid m\,]$, the probability that an earthquake of magnitude~$m$ produces surface rupture along the causative fault. Most implemented models in this category are empirical logistic regressions calibrated against historical surface-rupturing earthquakes \citep{WellsCoppersmith1994, youngs_methodology_2003, MossRoss2011, Petersen2011, Takao2013, Pizza2023}. The framework also implements the numerical approach of \citet{Mammarella2024} for all three faulting styles; this model is distinguished by an asterisk in Table~\ref{tab:model_library}. Numerical SR models complement empirical regressions in low-seismicity regions where the empirical record is sparse, and their inclusion in the logic tree allows the analyst to test the sensitivity of the hazard to the choice of approach. Figure~\ref{fig:principal_sr_probability} compares the implemented principal surface-rupture probability models across the supported faulting styles.


\begin{figure}[!htbp]
  \centering
  \includegraphics[width=\linewidth]{Figures/primary_surface_rupture_probability.pdf}
  \caption{Conditional probability of principal surface rupture versus moment magnitude $m$ for the implemented prediction models, organised by faulting style: (a) normal, (b) reverse, (c) strike-slip, and (d) style-independent (``all-slip''). Each curve is plotted only within its published applicable magnitude range \citep[Table~4]{Valentini2025}. Mammarella et al. (2024) curves use a seismogenic thickness of 9~km, dip angles of 50$^{\circ}$, 30$^{\circ}$, and 90$^{\circ}$ for normal, reverse, and strike-slip faulting respectively, and three magnitude-scaling relations.}
  \label{fig:principal_sr_probability}
\end{figure}

\subsection{Principal Fault Displacement}
\label{sec:model_library:principal_fd}

Principal displacement models predict $P[\,D > D_0 \mid m, x/L, sr \neq 0\,]$ given that surface rupture has occurred. This category is the most populated in the framework, reflecting the rapid expansion of model development through the FDHI Project \citep{Bozorgnia2021, Sarmiento2025}. Three groups can be distinguished. Pre-FDHI legacy models (\citealt{youngs_methodology_2003} for normal faulting, \citealt{Petersen2011} for strike-slip faulting, and \citealt{MossRoss2011} for reverse faulting) predict single-principal displacement components using continuous parametric distributions of normalised displacement. Recent FDHI-related models---\citet{LavrentiadisAbrahamson2023}, \citet{Kuehn2024}, \citet{Chiou2025}, and \citet{Moss2024}---predict net, sum-of-principal, or aggregate displacement using larger and more standardised data sets, and \citet{LavrentiadisAbrahamson2023} introduces the explicit gap and zero-displacement pathway described by Equation~\eqref{eq:zero_slip}. Regional models---including \citet{Takao2013} for Japanese reverse and strike-slip events---are implemented for completeness and for use in benchmark comparisons.

The variety of displacement metrics across these models means that selecting which models to combine in a logic tree is not arbitrary. The analyst must decide whether to mix legacy single-principal vertical or horizontal predictions with FDHI-related net or aggregate predictions, accept the resulting metric inconsistency as part of an explicitly stated epistemic spread, or restrict the tree to a single metric family. Figure~\ref{fig:principal_fdm_profiles} illustrates the corresponding heterogeneity in median principal displacement profiles for the implemented principal fault-displacement models.


\begin{figure}[!htbp]
  \centering
  \includegraphics[width=\linewidth]{Figures/principal_fdm.pdf}
  \caption{Median principal fault-displacement profiles versus normalised along-rupture position $x/L$, organised by faulting style. The top two rows show the normalised profiles $D/MD$ (a--c) and $D/AD$ (d--f); the bottom row shows absolute median displacement at $m=6.5$ on a logarithmic axis (g--i). Columns correspond to normal (a, d, g), reverse (b, e, h), and strike-slip (c, f, i) faulting. Profiles are symmetric about $x/L=0.5$, so only the half-profile is shown, and each curve is plotted only within its published applicable magnitude and $x/L$ range \citep[Table~4]{Valentini2025}. For Petersen et al. (2011) in panel~(f), the source-defined $D/AD$ lognormal coefficients are used; the model is absent from panel~(c) because the source publication does not define a $D/MD$ distribution. For Moss and Ross (2011), the Gamma and Weibull $D/AD$ variants are both plotted in panel~(e), and their Beta $D/MD$ variant is plotted in panel~(b). Moss et al. (2024) appears only in the normalised panels (b, e) because its source publishes only normalised $D/MD$ and $D/AD$ distributions.}
  \label{fig:principal_fdm_profiles}
\end{figure}


\subsection{Distributed Surface-Rupture Probability}
\label{sec:model_library:distributed_sr}

Distributed surface-rupture probability models predict $P[\,d \neq 0 \mid z, r, sr \neq 0\,]$, the probability that distributed rupture occurs within a site of dimension~$z$ at distance~$r$ from the principal fault. The framework implements the gridding-based formulation of \citet{FerrarioLivio2021} for normal faulting and Takao-family models for reverse and strike-slip faulting \citep{Takao2013, Takao2014}. The two Takao studies are retained as separate branches because they use different cell sizes and calibration data sets: \citet{Takao2013} is calibrated at 500~m, whereas \citet{Takao2014} is calibrated at both 500~m and 100~m. \citet{RodriguezPadillaOskin2023} provides an additional strike-slip distributed-rupture model based on historic earthquake ruptures along immature strike-slip faults. The current population of this category is sparse compared with the principal counterparts, reflecting the smaller pool of published distributed-only models; coupled distributed SR+FD pathways are described in Section~\ref{sec:model_library:distributed_srfd}. Figure~\ref{fig:distributed_sr_probability} compares the implemented distributed surface-rupture probability models---including the distributed surface-rupture probability component of the paired distributed SR+FD pathways of Section~\ref{sec:model_library:distributed_srfd}---across the supported faulting styles.


\begin{figure}[!htbp]
  \centering
  \includegraphics[width=\linewidth]{Figures/distributed_surf_rup.pdf}
  \caption{Probability of distributed (secondary) surface rupture versus distance from the principal fault trace $r$, organised by faulting style: (a) normal, (b) reverse, and (c) strike-slip. For dip-slip panels (a, b), $r$ is signed (negative = footwall, positive = hanging wall) and curves are drawn continuously across $r = 0$; the strike-slip panel (c) is shown for $r \geq 0$ only. Magnitude-dependent models are evaluated at $m = 6.5$. Panel (a) uses a $0$--$1$ probability axis; panels (b) and (c) use $0$--$0.5$. Each curve is plotted only within its published applicable magnitude and distance range \citep[Table~4]{Valentini2025}. Cell/pixel sizes are noted in the legend; in panel (a) all models are evaluated on a 500~m grid. \citet{Visini2025} is shown for all three combinations A, B, and C at the indicated pixel sizes. \citet{RodriguezPadillaOskin2023} is included in panel (c) even though it lies outside the \citet[Figure~12]{Valentini2025} inventory.}
  \label{fig:distributed_sr_probability}
\end{figure}


\subsection{Coupled Distributed Surface-Rupture and Displacement Models}
\label{sec:model_library:distributed_srfd}

The fourth row of Table~\ref{tab:model_library} groups studies that provide both the distributed surface-rupture probability and the distributed displacement model as a paired pathway. \citet{youngs_methodology_2003} for normal faulting and \citet{Petersen2011} for strike-slip faulting fall in this category through the gridding decomposition of Equation~\eqref{eq:distributed}, with the principal surface-rupture probability applied as a multiplicative gate. The slicing-based approach of \citet{Visini2025} is implemented across all three faulting styles through the dedicated \texttt{VisiniSecondaryCalculator} described in Section~\ref{sec:architecture:hazard_engine}; in this pathway the principal surface-rupture probability is not applied as a multiplicative factor because the regressions are calibrated exclusively on surface-rupturing earthquakes. \citet{Moss2022} provides a paired SR+FD set for reverse faulting derived from the GIRS-2022-05 study. The distributed surface-rupture probability components of these paired pathways are plotted alongside the stand-alone distributed SR models in Figure~\ref{fig:distributed_sr_probability}. Figure~\ref{fig:distributed_fd_profiles} compares the implemented distributed fault-displacement profiles across the supported faulting styles. Because the published distributed displacement models do not share a common displacement metric, the percentile of the plotted curve differs by model and is reported explicitly in the legend and caption; this heterogeneity is intrinsic to the source publications rather than a presentation choice.


\begin{figure}[!htbp]
  \centering
  \includegraphics[width=\linewidth]{Figures/distributed_surf_displ.pdf}
  \caption{Distributed (secondary) fault displacement versus distance from the principal fault trace $r$, organised by faulting style: (a) normal, (b) reverse, and (c) strike-slip. For dip-slip panels (a, b), $r$ is signed (negative = footwall, positive = hanging wall) and curves are drawn continuously across $r = 0$; the strike-slip panel (c) is shown for $r \geq 0$ only. Magnitude-dependent models are evaluated at $m = 6.5$, and the \citet{Visini2025} regression uses a default dip of 60$^{\circ}$ for both normal and reverse panels. The percentile of the plotted curve differs by model and is labelled explicitly in the legend: \citet{youngs_methodology_2003} is shown at the 85th and 95th percentiles of $d/MD$ per the source publication; \citet{Visini2025} is the median (50th percentile) of the lognormal $\ln(Y)$ regression for each combination A, B, and C; \citet{Moss2022} is the 50th-percentile $d/MD$ envelope (Table~5.7 of GIRS-2022-05) under both its gamma and envelope modes; and \citet{Petersen2011} is the median (50th percentile) recovered from the lognormal exceedance regression of Eq.~18 of the source publication. Each curve is plotted only within its published applicable magnitude and distance range \citep[Table~4]{Valentini2025}. Pixel size does not enter the \citet{Visini2025} median displacement regression---it modulates only the conditional rupture probability shown in Figure~\ref{fig:distributed_sr_probability}---so a single curve per combination is drawn; likewise, the \citet{Petersen2011} Eq.~18 regression contains no cell-size term, so panel (c) shows a single curve. The \citet{Visini2025} curves spike at $r \approx 0$ because the ln$(s)$ regression term diverges at the principal trace and the TPFm smoothing window collapses to a point at the same limit; both are intrinsic to the published model rather than numerical artifacts. \citet{Moss2022} is included in panel (b) although it lies outside the \citet[Figure~16]{Valentini2025} inventory.}
  \label{fig:distributed_fd_profiles}
\end{figure}


\subsection{Extending the Library}
\label{sec:model_library:extending}

Adding a new model to the framework requires three steps: (i)~implement a Python class that subclasses the appropriate base class for one of the four PFDHA components and provides a \texttt{get\_prob()} method returning the relevant probability; (ii)~make the class available to the model lookup mechanism under a unique name; and (iii)~reference the class by name from a \texttt{logicTreeBranch} in the NRML logic-tree file, supplying any per-branch parameters in the inline TOML-style block described in Section~\ref{sec:architecture:logic_tree}. No modification to the calculation engine is required, and no existing models are affected. This extension path is intended to lower the barrier to community contributions, including reimplementation of regional models, future calibration or implementation of models based on the FDHI and SURE~2.0 databases, and inclusion of forthcoming non-ergodic formulations \citep{Abrahamson2024}.


% =====================================================================
% TABLE 1  (tabularx layout â fits page width automatically)
% =====================================================================

% =====================================================================
% TABLE 1 â robust SRL/SSA-compatible table block
% Uses a normal table float, a manual caption, and tabularx.
% This avoids the SSA-SRL side-caption formatting that can wrap long
% captions vertically and avoids the \tbl + tabularx interaction that can
% suppress the table body in Prism.
% Required packages in 00_Main.tex: tabularx, booktabs, array.
% =====================================================================
\begin{table}[!htbp]
  \centering
  \scriptsize
  \renewcommand{\arraystretch}{1.14}
  \setlength{\tabcolsep}{3pt}

  \refstepcounter{table}\label{tab:model_library}
  \begin{minipage}{\linewidth}
  \raggedright
  {\bfseries Table~\thetable.} Implemented PFDHA model library organised by component and faulting style.\par\vspace{0.6em}
  \end{minipage}

  \begin{tabularx}{\linewidth}{@{}>{\raggedright\arraybackslash}p{2.0cm} >{\raggedright\arraybackslash}X >{\raggedright\arraybackslash}X >{\raggedright\arraybackslash}X@{}}
    \toprule
    \textbf{Component} & \textbf{Normal} & \textbf{Reverse} & \textbf{Strike-slip} \\
    \midrule

    \textit{Principal SR}\newline
    {\scriptsize $P[\,sr \neq 0 \mid m\,]$} &
      Wells and Coppersmith, 1994\newline
      Youngs et al., 2003\newline
      Pizza et al., 2023\newline
      Mammarella et al., 2024$^{*}$ &
      Moss and Ross, 2011; Moss et al., 2013\newline
      Takao et al., 2013\newline
      Yang et al., 2021\newline
      Mammarella et al., 2024$^{*}$ &
      Petersen et al., 2011\newline
      Moss and Ross, 2011; Moss et al., 2013\newline
      Takao et al., 2013\newline
      Mammarella et al., 2024$^{*}$ \\

    \midrule

    \textit{Principal FD}\newline
    {\scriptsize $P[\,D > D_0 \mid \cdot\,]$} &
      Youngs et al., 2003$^{v}$ (S)\newline
      Lavrentiadis and Abrahamson, 2023$^{n}$ (A)\newline
      Kuehn et al., 2024$^{n}$ (A) &
      Takao et al., 2013$^{v}$ (S)\newline
      Moss et al., 2022$^{v}$ (S)\newline
      Moss et al., 2024$^{v}$ (S)\newline
      Lavrentiadis and Abrahamson, 2023$^{n}$ (A)\newline
      Kuehn et al., 2024$^{n}$ (A) &
      Petersen et al., 2011$^{h}$ (S)\newline
      Chiou et al., 2025$^{n}$ (A)\newline
      Lavrentiadis and Abrahamson, 2023$^{n}$ (A)\newline
      Kuehn et al., 2024$^{n}$ (A) \\

    \midrule

    \textit{Distributed SR}\newline
    {\scriptsize $P[\,d \neq 0 \mid z, r, sr \neq 0\,]$} &
      Ferrario and Livio, 2021 &
      Takao et al., 2013\newline
      Takao et al., 2014 &
      Takao et al., 2013\newline
      Takao et al., 2014\newline
      Rodriguez Padilla and Oskin, 2023 \\

    \midrule

    \textit{Distributed}\newline \textit{SR+FD paired} &
      Youngs et al., 2003$^{v}$\newline
      Visini et al., 2025$^{n}$ &
      Visini et al., 2025$^{v}$\newline
      Moss et al., 2022$^{v}$ &
      Petersen et al., 2011$^{h}$\newline
      Visini et al., 2025$^{v}$ \\

    \bottomrule
  \end{tabularx}

  \vspace{0.55em}
  \begin{minipage}{\linewidth}
  \raggedright\scriptsize
  \textit{Notes.} Superscripts indicate displacement component:
  $^{v}$ vertical, $^{h}$ horizontal (lateral), and $^{n}$ net.
  Parenthetical labels indicate the displacement-metric family following
  Sarmiento et al., 2025: (S) single-principal and
  (A) sum-of-principal or aggregate. The asterisk ($^{*}$) marks numerical
  or physics-based models.
  \end{minipage}
\end{table}
